\section{Open thermodynamic and AD derivative questions}
\label{sec:open_questions}

This section records issues that must be resolved before declaring the MOOSE
implementation thermodynamically faithful.

The automatic-differentiation strategy is no longer an open yes-or-no question:
MOOSE can use \texttt{ADDerivativeParsedMaterial} for scalar free-energy
derivatives, while tensor kinematics, phase closure algorithms, flux laws,
constraints, and weak forms still require explicit AD-aware implementation.
The open questions below specify where that division must be made precise.

\begin{enumerate}[leftmargin=*]
	\item \textbf{Primitive-variable Jacobian.}  Determine which primitive
	variables give robust AD derivatives for \(\bar p_f\),
	\(\mu_\xi^\alpha\), \(\bar p_E\), \(B\), and the phase fractions while
	respecting the volume and mass-fraction constraints.  This is open for the
	full theory, but the first closed problem uses solid-reference component
	storage variables, partition variables, and conversion potentials until a
	better primitive set is justified.

	\item \textbf{Held-fixed variables.}  For every derivative inherited from the
	Helmholtz free energy, record what is held fixed: entropy or temperature,
	component masses or mass fractions, intrinsic densities, solid distension,
	and skeleton deformation.  These choices control the stress and pressure
	terms in the Newton matrix.

	\item \textbf{Equivalent pressure and nonlinear Biot coefficient.}  Verify
	the AD path for \(\bar p_E\) and \(B\) against the theory manuscript,
	including the distinction between pore-volume filling and stress-generating
	skeleton deformation.

	\item \textbf{Reaction-rate derivatives.}  Decide whether kinetics depend on
	current volume fractions, Lagrangian component masses, chemical potentials,
	affinities, or activities.  The implementation should include the full AD
	derivative of the chosen closure.

	\item \textbf{Flux derivative completeness.}  Check derivatives of
	\(\mbf W_f\) with respect to \(\mbf F\), \(J\), phase fractions, intrinsic
	densities, potentials, pressure-like variables, reaction rates, and
	permeability.  The pulled-back permeability
	\(J\mbf F^{-1}\widetilde{\bs\kappa}_f\mbf F^{-T}\) is a likely source of
	missing terms, especially because
	\(\widetilde{\bs\kappa}_f\) itself contains
	\(\sum_\alpha\dot c_f^\alpha\mbf I\).

	\item \textbf{Fluid phase kinematics.}  Decide whether any later closure
	requires \(J_f\) or \(\mbf F_f\) as a finite-element variable, material
	internal variable, or user-object state.  The first review problem does not
	require those fields as residual unknowns; it discretizes solid-frame
	advection directly by \(\mbf W_f/(J\rho_f)\).  Separate \(J_f\) or
	\(\mbf F_f\) evolution should be added only if a later constitutive closure
	needs the full fluid phase map or phase-reference history variables.

	\item \textbf{Constraint enforcement.}  Compare variable elimination,
	multiplier residuals, and constrained material solves for the volume-fraction
	and mass-fraction constraints.  The selected method must preserve the
	thermodynamic interpretation of pressure and component potentials.

	\item \textbf{Consistent units.}  Audit whether each source term is per
	current mixture volume or per solid reference volume before multiplying by
	\(J\).  This is especially important for reaction rates, conversion
	momentum, and boundary fluxes.
\end{enumerate}

These questions should remain visible in the paper until each is answered by a
derivation, a source-code trace, and at least one verification case.
